\section{Complete Method Configuration Matrix}
\label{app:method-matrix}

This appendix records the configuration contract used by the confirmation protocol. All methods receive the same prepared stream, four-window warm-up, one-window label delay, and the same model and stream seeds. The audit identified method-specific supervision semantics: all locked paths record 128 new training examples per update, but primary DASP constructs that subset from a labeled repair pool and additionally evaluates a disjoint 25\% labeled performance holdout; primary SSF also selects from a delayed labeled pool, while several comparator paths use label-blind uniform selection. Therefore, the current matrix does not establish a strict shared total-label-exposure contract.

\begin{table*}[t]
\caption{External evaluation controls and the audited DASP supervision exception.}
\label{tab:shared-contract}
\centering
\begin{tabular}{p{0.38\textwidth}p{0.18\textwidth}p{0.35\textwidth}}
\toprule
Control & Value & Rationale\\
\midrule
Task & Binary benign/attack & Common label space for all methods\\
Input features & 78 for CIC-IDS2017; 49 for NF-v3 streams & Dataset-specific fixed representation interface\\
Window size & 4,096 flows & Identical temporal granularity\\
Warm-up & 4 windows & Initial reference and classifier calibration\\
Label delay & 1 window & Prediction-first deployment constraint\\
Recorded new training examples & 128/window & Common optimization-subset count, not common total label exposure\\
DASP repair-training subset & At most 128/window & Class-balanced selection from the labeled repair pool\\
Stream seeds & 2001, 2002, 2003, 2101, 2102 & Replication over stream construction\\
Model seeds & 91, 92 & Replication over initialization\\
Device & CPU, one thread & Deterministic performance measurement\\
Memory capacity & 1,024 records & Bounded historical state\\
DASP performance holdout & 0.25 of delayed window & Labeled pre-update emergency/background mode selection; outside the 128-example training subset\\
Primary metric & Macro-F1 & Balances benign and attack classes\\
\bottomrule
\end{tabular}
\end{table*}

\begin{table*}[t]
\caption{Method-specific controls used in the confirmation matrix.}
\label{tab:method-config-detail}
\centering
\scriptsize
\begin{tabular}{p{0.18\textwidth}p{0.22\textwidth}p{0.22\textwidth}p{0.25\textwidth}}
\toprule
Method & Adaptation unit & Historical protection & Distinguishing control\\
\midrule
DASP & Attributed modules plus classifier & Class-balanced 512-example replay & MMD/conflict router, top-$k=3$, strict optimizer exclusion\\
SSF & Full predictor update & Strategic sample selection and forgetting protection & Fixed sample mask, distillation/forgetting control\\
Experience Replay & Full predictor update & Replay memory & Historical examples mixed with current labels\\
RepShield & Representation-protection path & Protected representation reference & Robust representation preservation\\
A-GEM & Full predictor update & Reference gradient memory & Gradient projection against memory direction\\
Online MLP & Full predictor update & None beyond current state & Direct online supervised update\\
Online EWC & Full predictor update & Importance-weighted regularization & EWC penalty on important parameters\\
CDDA-MD & Detector-specific update & Distributional reference & MMD-style drift detector with compact update\\
\bottomrule
\end{tabular}
\end{table*}

The matrix clarifies a potential source of confusion in comparative studies. DASP's trainable unit is not equivalent to using a smaller batch. DASP still evaluates the full modular predictor at prediction time and still uses historical examples during a repair. The distinction is where gradients are allowed to change parameters. Similarly, a baseline with fewer gradient evaluations is not automatically more efficient in the operational sense if it obtains substantially lower recall or performs almost no adaptation.

\subsection{Locked DASP values}
The final DASP configuration uses four modules, module width 64, hidden dimensions 128 and 64 in the surrounding implementation, no dropout, batch size 128, AdamW learning rate $10^{-3}$, weight decay $10^{-4}$, and two triggered/background repair epochs. The router samples 128 module representations, uses four calibration windows, kernel scales $\{0.5,1,2\}$, severity threshold 0.15, persistence one, cooldown zero, and top-$k=3$. The emergency-severity value 0.60 in the locked experimental configuration is inactive in the evaluated primary procedure. Each repair constructs a class-balanced optimization subset of at most 128 new examples from the labeled repair pool, evaluates a disjoint 25\% labeled performance holdout, and samples 512 replay examples. Background active modules use plasticity 0.25; emergency modules use $\min\{1,\max(0.25,c_j)\}$ for three epochs. The classifier multiplier is 1 on the base learning rate. Reference refresh is enabled after every applied primary DASP repair.

These values should be regarded as part of the method definition for the reported result. Changing the memory size, router budget, or top-$k$ changes the method and requires a new confirmation experiment. In particular, a full-model update should not be described as DASP merely because the update began with a module-wise score.

\section{Dataset and Scenario Design}
\label{app:scenario-manifest}

The following design summary explains what is varied and what is held fixed in the prepared streams. A scenario is a transformation of an ordered stream, not a random label permutation. The transformation is applied before model training and is identical for all methods within a paired cell.

\begin{table*}[t]
\caption{Dataset characteristics and intended stress properties.}
\label{tab:dataset-manifest}
\centering
\begin{tabular}{p{0.25\textwidth}p{0.28\textwidth}p{0.32\textwidth}}
\toprule
Dataset & Stream role & Main stress property\\
\midrule
CIC-IDS2017 & Enterprise benchmark & Multiple attack families and benign variability\\
NF-BoT-IoT-v3 & Botnet/DoS NetFlow benchmark & Localized botnet and denial-of-service changes\\
NF-CSE-CIC-IDS2018-v3 & Enterprise/IoT NetFlow variant & Heterogeneity with high aggregate separability\\
NF-ToN-IoT-v3 & Telemetry-rich IoT benchmark & Difficult evolving traffic composition\\
NF-UNSW-NB15-v3 & Heterogeneous modern-flow benchmark & Ceiling performance and small residual gaps\\
\bottomrule
\end{tabular}
\end{table*}

\begin{table*}[t]
\caption{Scenario design with temporal behavior, expected failure mode, and DASP diagnostic.}
\label{tab:scenario-manifest-detail}
\centering
\scriptsize
\begin{tabular}{p{0.16\textwidth}p{0.29\textwidth}p{0.25\textwidth}p{0.23\textwidth}}
\toprule
Scenario & Temporal behavior & Expected challenge & Diagnostic emphasis\\
\midrule
S0 & No deliberate regime change & Unnecessary plasticity & Background-mode stability\\
S1-Sudden & Abrupt transition & Detection latency & Immediate attribution\\
S2-Gradual & Distributed transition & Ambiguous attribution & Score persistence and worst cell\\
S3-Incremental & Several smaller changes & Repeated repairs & Replay and cumulative forgetting\\
S4-Recurrent & Return of a prior regime & Overwriting and recovery & Replay-conflict behavior\\
S5-Blip & Brief transient regime & Overreaction & Emergency/background mode behavior\\
S6-Combined & Compound transformation & Interacting drift signals & Module-count budget\\
S7-Local-Low & Weak local change & Low signal-to-noise ratio & Positive sparse repair\\
S7-Local-Medium & Moderate local change & Boundary movement & Attribution precision\\
S7-Local-High & Strong local change & Multiple-module activation & Top-$k$ limitation\\
\bottomrule
\end{tabular}
\end{table*}

The scenario labels are used in the analysis as scientific conditions, not as method-specific hints. DASP does not receive the string ``S7-Local-Low'' as an input. It receives only the stream observations and delayed labels. The scenario name is available to the experiment manager for grouping after the run completes.

\subsection{Protocol validation}
Before a run is accepted, the evaluation procedure verifies that the dataset, scenario, stream seed, and feature count agree with the requested experimental cell. It also verifies that warm-up windows precede the first adaptation window. A mismatch produces an invalid run rather than a silently relabeled result. This is particularly important in a 4,000-run study, where an indexing error could otherwise appear as a statistical effect.

\subsection{Class balance and label accounting}
The memory sampler maintains class balance whenever both classes are available, and pseudo-labels do not substitute for delayed ground truth. On the primary DASP path, class balancing inspects repair-pool labels before selecting at most 128 new optimization examples, and a disjoint 25\% labeled performance holdout is evaluated separately. The recorded 128 field is therefore not a total-label-exposure count. Replay examples remain historical memory rather than new supervision. A strict supervision-matched confirmation requires a rerun in which every method follows one common label-access and holdout contract.

\section{Prediction-First State Machine}
\label{app:state-machine}

The prediction-first contract is indexed by the outer iteration $t$ and delayed repair index $\tau=t-1$. The model first predicts $W_t$, computes representation evidence $\mathcal E_t$, and enqueues the current item. It then receives $Y_{\tau}$, restores the queued evidence $\mathcal E_{\tau}$ associated with the same delayed window, computes module-wise gradient conflict, and selects triggered attributed repair, performance-gated emergency repair, or background local consolidation. This textual state machine replaces the older auxiliary diagram so that Figure~\ref{fig:window-algorithm} remains the single canonical algorithm figure.
\subsection{State transition conditions}
The current-window prediction and evidence computation are unconditional and occur before access to $Y_t$. The delayed transition is controlled by the fixed one-window queue. Triggered attributed repair is controlled by the stored MMD decision $D_{\tau}$; within that path, conflict attenuates the stored requirement. When $D_{\tau}$ does not trigger, delayed-label performance-holdout Macro-F1 selects emergency or background local repair.

After the routed update, the repair-training subset is added to bounded memory, the enabled reference state is refreshed from the applied model and $W_{\tau}$, and the routing record is stored. The evaluated DASP procedure does not implement post-update candidate rejection or rollback.

\subsection{Expanded pseudocode}
\begin{algorithmic}[1]
\STATE Initialize $f_{\Theta}$, $\mathcal H^{\mathrm{ref}}$, memory $\mathcal M$, queue $\mathcal Q$, and audit log.
\FOR{outer iteration $t=1,\ldots,T$ with current window $W_t$}
  \STATE $\hat Y_t\leftarrow f_{\Theta_t}(X_t)$; emit predictions before label access.
  \STATE Compute $\mathcal E_t=(\mathbf r_t,s_t,D_t)$ using the current window and reference snapshot.
  \STATE Push $(W_t,\hat Y_t,\mathcal E_t)$ into $\mathcal Q$.
  \IF{labels are available for $\tau=t-1$}
    \STATE Pop $(W_{\tau},\hat Y_{\tau},\mathcal E_{\tau})$ and restore $\mathbf r_{\tau},s_{\tau},D_{\tau}$.
    \STATE Split the delayed window into repair pool and disjoint performance holdout.
    \STATE Use repair-pool labels to construct class-balanced $C_{\tau}^{\mathrm{new}}$ with at most 128 training examples; draw $B_{\tau}^{\mathrm{mem}}$.
    \STATE Compute $c_{j,\tau}$ from $C_{\tau}^{\mathrm{new}}$ and $B_{\tau}^{\mathrm{mem}}$ gradients.
    \IF{stored MMD decision $D_{\tau}$ triggers}
      \STATE Set $p_{j,\tau}=r_{j,\tau}/(1+c_{j,\tau})$ and route positive top-$k$ modules.
    \ELSIF{$F^{\mathrm{hold}}_{\tau}<0.80$}
      \STATE Route top-$k$ conflicts for performance-gated emergency repair.
    \ELSE
      \STATE Route top-$k$ conflicts for background local consolidation.
    \ENDIF
    \STATE Construct $\widetilde C_{\tau}=C_{\tau}^{\mathrm{new}}\cup B_{\tau}^{\mathrm{mem}}$ with new-sample weight 3.
    \STATE Exclude all $\Theta_j$, $j\notin\mathcal A_{\tau}$, from AdamW; train active modules and classifier.
    \STATE Update memory and audit; refresh the enabled reference state using $W_{\tau}$.
    \STATE Store the updated model for subsequent windows; no candidate rollback is executed.
  \ENDIF
\ENDFOR
\end{algorithmic}

The pseudocode makes two implementation details explicit. First, delayed-label access occurs after prediction, not merely after feature extraction; the 128-example value bounds the optimization subset rather than total observed labels. Second, the optimizer parameter list changes per repair. A post-hoc mask applied to gradients would not by itself guarantee protection from optimizer state or weight decay.

\section{Statistical Audit of Paired Results}
\label{app:statistical-audit}

The primary comparison contains 500 paired DASP--SSF cells. The distribution of paired Macro-F1 gaps is asymmetric but centered above zero. The minimum gap is approximately $-0.0900$, the first quartile is approximately $+0.0002$, the median is $+0.0032$, the third quartile is approximately $+0.0131$, and the maximum is approximately $+0.0684$. The quartiles show that the positive mean is not caused only by a single extreme win, while the negative minimum confirms that individual failures remain operationally relevant.

\begin{table*}[t]
\caption{Distribution audit for the 500 paired DASP--SSF Macro-F1 gaps.}
\label{tab:gap-quantiles}
\centering
\begin{tabular}{lrrrrrr}
\toprule
Statistic & Minimum & Q1 & Median & Q3 & Maximum & Mean\\
\midrule
Macro-F1 gap & -0.0900 & 0.0002 & 0.0032 & 0.0131 & 0.0684 & 0.0078\\
\bottomrule
\end{tabular}
\end{table*}

\begin{table*}[t]
\caption{Paired adaptation-cost distribution across the same 500 cells.}
\label{tab:cost-quantiles}
\centering
\begin{tabular}{lrrrrrr}
\toprule
Metric and method & Minimum & Q1 & Median & Q3 & Maximum & Mean\\
\midrule
DASP update seconds & 1.248 & 1.319 & 1.357 & 1.404 & 5.750 & 1.415\\
SSF update seconds & 2.428 & 2.514 & 2.591 & 2.680 & 9.229 & 2.683\\
DASP gradient evaluations & 320 & 320 & 320 & 320 & 410 & 328.7\\
SSF gradient evaluations & 512 & 512 & 512 & 512 & 512 & 512.0\\
\bottomrule
\end{tabular}
\end{table*}

The gradient-evaluation distribution is especially interpretable. SSF uses 512 evaluations in each paired cell under the locked protocol, whereas DASP is usually at 320 and occasionally increases to 410 when the repair path requires additional work. The cost reduction is therefore not an artifact of a few unusually fast cells. It follows from the module mask and the bounded repair schedule.

\subsection{Multiple-comparison discipline}
The sign test, paired $t$-test, and Wilcoxon test are used as complementary analyses of the same primary paired comparison. They should not be presented as three independent discoveries. The sign test is the most robust headline test because it uses only the direction of each paired difference. The other tests provide sensitivity to magnitude and rank structure. Exploratory dataset- and scenario-level comparisons are descriptive and should be interpreted with Holm correction if formal claims are made across many groups.

\subsection{What is not treated as independent}
The 16,000 window-level repair records are not counted as 16,000 independent hypothesis tests. They are used to verify structural invariants and diagnose routing behavior. Likewise, the 4,000 runs are not treated as 4,000 independent samples for the DASP--SSF superiority claim; the paired claim uses the 500 cells containing both methods. This distinction prevents pseudo-replication and keeps the inferential unit aligned with the experimental design.

\section{Performance and Memory Audit}
\label{app:performance-audit}

The performance results have three distinct layers. Prediction throughput measures the rate at which a trained model can score flows. Adaptation time measures the cost of processing delayed labels and updating the state. Memory and resident-set measurements describe the operating footprint. A method can be favorable on one layer and unfavorable on another; therefore no single scalar efficiency score is reported.

\begin{table*}[t]
\caption{Confirmation-level performance summary by method.}
\label{tab:performance-audit}
\centering
\scriptsize
\begin{tabular}{lrrrrrr}
\toprule
Method & Macro-F1 & Update s & Throughput & Grad. eval. & Memory payload & Peak RSS increment\\
\midrule
DASP & 0.9269 & 1.415 & 159,219 & 328.7 & 320 KB & 28.6 MB\\
SSF & 0.9191 & 2.683 & 455,279 & 512.0 & 320 KB & 28.9 MB\\
Experience Replay & 0.9036 & 0.546 & 461,493 & 192.0 & 320 KB & 28.4 MB\\
RepShield & 0.8881 & 0.946 & 464,789 & 128.0 & 320 KB & 30.1 MB\\
A-GEM & 0.8821 & 0.495 & 464,700 & 128.0 & 320 KB & 29.1 MB\\
Online MLP & 0.8751 & 0.222 & 464,101 & 64.0 & 320 KB & 31.1 MB\\
Online EWC & 0.8751 & 0.396 & 454,905 & 64.0 & 320 KB & 28.9 MB\\
CDDA-MD & 0.7973 & 0.038 & 458,585 & 2.7 & 320 KB & 27.9 MB\\
\bottomrule
\end{tabular}
\end{table*}

DASP's throughput penalty is an architectural consequence of computing four module branches and a fusion classifier. It should not be described as a memory penalty: the model parameter footprint and memory payload are comparable to the other methods in the locked implementation. The practical implication is that a deployment must choose whether the critical resource is inference latency or adaptation latency. A future optimized implementation could fuse the four branches for prediction while retaining separate parameter groups for repair, or could use sparse module execution when the router predicts that some groups are inactive for a known traffic class.

\subsection{Update-time interpretation}
The mean update time includes all update intervals. This includes discrepancy calculation, replay sampling, localized training, performance-holdout evaluation on the no-trigger path, and any emergency path. Reporting the mean alone can hide the effect of mode-dependent work, so the median and maximum are included in Table~\ref{tab:cost-quantiles}. The maximum DASP update time is larger than its mean because a small number of difficult cells invoke extra work. The same pattern is present for SSF, but its entire distribution is shifted upward under the locked schedule.

\subsection{Memory interpretation}
The replay capacity is 1,024 records, while a DASP repair uses 512 records. The memory payload reported in the summary is therefore bounded and does not grow with stream length. Peak RSS is affected by the runtime, tensor allocations, and input-data loading, so it is interpreted as an engineering diagnostic rather than a pure model-size measure. The manuscript reports it to make the resource trade-off visible, not to claim a universal memory advantage.

\section{Reproduction and Archival Guide}
\label{app:archival-guide}

A complete reproduction requires four components: chronologically prepared streams, the locked experimental configuration, executable method implementations, and row-level plus aggregate analysis records. The archival release should provide these components together with instructions for reproducing the confirmation matrix, statistical comparisons, tables, and figures.

\begin{table*}[t]
\caption{Archival artifact map.}
\label{tab:artifact-map}
\centering
\begin{tabular}{p{0.36\textwidth}p{0.28\textwidth}p{0.25\textwidth}}
\toprule
Artifact & Function & Validation use\\
\midrule
Chronologically prepared streams & Ordered windows and labels & Verify chronology and feature count\\
Stream metadata records & Seed and scenario metadata & Verify pairing and experimental-cell identity\\
Locked experimental configuration & Fixed hyperparameters & Verify method contract\\
Run summaries & One row per method/cell & Recompute aggregates\\
Paired cell-level comparison & DASP--SSF join & Recompute tests and quantiles\\
Protocol-level confirmation record & Run-count and invariant audit & Verify 4,000-run matrix\\
Generated figures & Visual result summaries & Verify plotted aggregates\\
LaTeX source & Manuscript and equations & Recompile article\\
\bottomrule
\end{tabular}
\end{table*}

\subsection{Recommended execution order}
The archival sequence should begin with a small-scale evaluation on one dataset, one scenario, one stream seed, one model seed, and all methods. This initial evaluation verifies software availability, configuration resolution, prediction ordering, run-record structure, and basic metric calculation. The full pilot then covers five datasets, ten scenarios, three stream seeds, one model seed, and eight methods. Only after the full pilot is complete should the confirmation be executed with five stream seeds and two model seeds. The confirmation records are then frozen before manuscript tables are regenerated.

\subsection{Failure handling}
A failed run should be marked invalid and repeated from the same clean seed rather than silently replaced by a partial summary. If a machine-level failure occurs, the evaluation procedure should retain a diagnostic record and the resolved configuration. A result should not enter the paired analysis unless both DASP and SSF completed the same cell. This rule is essential for preserving the meaning of the win count.

\subsection{Reporting checklist}
Before submission, the authors should replace metadata placeholders, verify all dataset licenses, archive the exact software environment, record the CPU model, and confirm that the final PDF uses the same finalized result records as the source tables. The manuscript should also retain the negative cell and the throughput limitation. Removing either would make the narrative less credible, even if the page count became shorter.

\section{Extended Limitations and Future Experiments}
\label{app:future-work}

The confirmation establishes a useful but bounded result. The streams are controlled chronological artifacts rather than a live multi-site deployment. The feature partition is fixed and contiguous, so the modules are operational subspaces rather than guaranteed semantic concepts. The binary label space does not test retention of individual attack families. The model seeds and stream seeds provide replication but do not exhaust preprocessing, architecture, or hardware variability. These limitations define the next experiments rather than being reasons to reinterpret the current result.

\begin{table*}[t]
\caption{Prioritized future experiment matrix.}
\label{tab:future-experiments}
\centering
\scriptsize
\begin{tabular}{p{0.24\textwidth}p{0.27\textwidth}p{0.25\textwidth}p{0.14\textwidth}}
\toprule
Experiment & Manipulation & Primary outcome & Priority\\
\midrule
Adaptive $k$ & $k\in\{1,2,3,4\}$ by uncertainty & Accuracy versus active-module fraction & High\\
Learned grouping & Replace contiguous groups with learned groups & Attribution stability & High\\
Causal router ablation & Shuffle or mask router scores & Accuracy loss and locality & High\\
Natural streams & Replace controlled scenarios with chronological sites & External validity & High\\
Multi-class retention & Attack-family labels and open-set attacks & Class-wise forgetting & High\\
Sparse inference & Skip inactive prediction branches & Throughput recovery & Medium\\
Label noise & Corrupt delayed labels at several rates & Performance-gate robustness & Medium\\
Hardware scaling & CPU/GPU and batch-size sweep & Systems portability & Medium\\
\bottomrule
\end{tabular}
\end{table*}

The causal router ablation is particularly important for a journal extension. The current evidence shows that the implementation is structurally local and that the final method performs well. It does not isolate the contribution of the representation discrepancy, supervised conflict, or top-$k$ mask separately. A causal ablation would preserve the same replay, performance-gate, top-$k$, label, and optimization budgets while comparing random top-$k$, fixed-module, MMD-only, conflict-only, and combined attribution routing. An all-module variant should be reported separately as a locality ablation because it removes the top-$k$ structural constraint. Such an experiment would turn the attribution narrative from a plausible mechanism into a directly tested causal claim.

The natural-stream experiment is equally important for deployment. Controlled scenarios provide clean diagnostic contrast, but operators encounter changes in service mix, routing policy, benign workload, and attack tooling that do not fit a single synthetic taxonomy. A natural-stream study should keep the prediction-first and delayed-label rules, report site-level pairing, and treat collection site as a higher-level replication factor. It should not simply concatenate all data and call the result ``realistic.''

Finally, the throughput limitation should be treated as an engineering target. The current architecture executes all branches to obtain attribution-ready representations. A production version could maintain a full-resolution path during adaptation and a fused inference path between repairs. The correctness condition would be that fusion does not change the logits used for the reported predictions, or that any approximation is measured as a separate deployment variant.


